\medskip\noindent
\textbf{The actual homotopy group and a stable cell model.}
The cover used above has an actual double-loop realization.
Put $Z=\operatorname{fib}(Y\to K(\mathbb F _2,4))$.
Then $F=\Omega^2Z$, the space $Z$ is five-connected, and
$\pi_6Z=\ZZ/2$.  Double-looping its first Postnikov map
$Z\to K(\mathbb F _2,6)$ realizes the map
$F\to K(\mathbb F _2,4)$, whose fiber is $E$.
Thus $E$ is a four-connected double-loop space.
The unit $E\to\Omega\Sigma E$ has a homotopy retraction by
the loop-space counit; Freudenthal makes it nine-connected.
Since $\Sigma E$ is five-connected, this proves
\[
 \pi_9T\cong\pi_9F\cong\pi_9E
       \cong\pi_9\Sigma^\infty E.
 \tag{D.2.22t1}
\]
This argument does not assert that stabilization of $F$ preserves
its specified double-loop operations.

Work two-locally.  Choose a finite stable approximation to
$\Sigma^\infty E$ through dimension eleven before cancelling
odd-unit cell pairs, and retain its ten-skeleton.
Write $L_i$ for its stable skeleta.  This convention removes
extraneous ten-cells killed by unit eleven-cell incidences and
preserves the required $\pi_9$.
Let $M=M(\ZZ/4,5)$, with bottom inclusion $i_5$ and top
six-class $D$, and put $r=D+E_6$ and $t=i_5\eta_5$.
The lower integral boundaries and operations in the preceding
tables give
\[
 \begin{split}
 D_0&=(M\vee S^6_r)\mathop{\cup}_{t+2r}e^7_S,
       \qquad L_7=D_0\vee S^7_R,\\
 D_0&=\cofib\left(
 \begin{pmatrix}4&\eta\\0&2\end{pmatrix}:
 S^5\vee S^6\longrightarrow S^5\vee S^6\right).
 \end{split}
 \tag{D.2.22t2}
\]
In particular $\beta S=r$ and $\operatorname{Sq}_*^2S=u_5$.
All cell assertions here concern stable cofibers; no minimal
unstable CW decomposition is being assumed.

Let $b\in\pi_7M$ be the order-two lift of $\eta$ on the
Moore top, and let $c=r\eta_6$.
The cofiber exact sequence gives
\[
 \pi_7L_7=\ZZ_{(2)}\{R\}\oplus
                    \ZZ/2\{b,c\}.
\]
The complete eight-cell attachment list, ordered as
$W,D_8,J_1,J_2,J_3$, is
\[
 0,\qquad b+c,\qquad b,\qquad0,\qquad2R+b.
 \tag{D.2.22t3}
\]
The $b,c$ coordinates are detected independently by the
second-square vectors $D,r$; the integral coordinate is fixed
by $\beta J_3=R$.
Subtract the $J_1$ cell from each of $D_8,J_3$.
These stable cofiber row changes give an actual splitting
\[
 \begin{split}
 L_8&=K_8\vee M(\ZZ/2,7)_R\vee S^8_W\vee S^8_{J_2},\\
 K_8&=D_0\cup_b e^8_{B_8}\cup_c e^8_{C_8},
       \qquad B_8=J_1,\quad C_8=D_8-J_1.
 \end{split}
 \tag{D.2.22t4}
\]
The subscripts $B_8,C_8$ here distinguish these cells from
the original space $B$ and from the cover's ten-classes $B_i$.

\medskip\noindent
\textbf{All ninth homotopy components of the eight-skeleton.}
First, $\pi_9M=\ZZ/4\{m\}$, with top quotient
$q_*m=2\nu_6$.  Write $z=r\nu_6$.
The full matrix cofiber sequence identifies the quotient of
$\pi_9D_0$ by $\ZZ/2\{z\}$ with the kernel of
\[
 \ZZ/8\{\nu_5\}\oplus\ZZ/2\{\eta_6^2\}
       \longrightarrow\ZZ/8,
       \qquad (x,y)\longmapsto4x+4y.
 \tag{D.2.22t5}
\]
This kernel is cyclic of order eight generated by
$(\nu_5,\eta_6^2)$; choose a lift $s_0$.
The possible discrepancy in $2s_0=m$ is $z$.
It is zero: the fourth-square detector of both $2s_0$ and
$m$ is zero, whereas the detector of $z$ is $r$.
The same argument excludes $8s_0=z$.
Changing $s_0$ by $z$ if necessary normalizes its
fourth-square vector to $D$.  Consequently
\[
 \pi_9D_0=\ZZ/8\{s_0\}\oplus\ZZ/2\{z\},
       \qquad2s_0=m.
 \tag{D.2.22t6}
\]
Here a square detector means the homology square on the top
class of the stable mapping cone.  In this range it is additive
and vanishes on twice an attaching map: the map from the cone
of $2f$ to the cone of $f$ is zero on the top mod-two class
and is the identity on the relevant lower homology.
The nonzero detector for $z$ is the Hopf-invariant-one
detector of $\nu$.

The two eight-attachments have composites
\[
 b\eta_7^2=2m=4s_0,
       \qquad c\eta_7^2=4z=0.
\]
There is no additional $\eta$-lift in their cofiber sequence:
the classes $b\eta_7,c\eta_7$ are independent before those
eight-cells are attached.  Therefore
\[
 \begin{split}
 \pi_9K_8&=\ZZ/4\{s\}\oplus\ZZ/2\{z\},\\
 \pi_9L_8&=\ZZ/4\{s\}\oplus\ZZ/2\{z\}
     \oplus\ZZ/4\{\tau\}
     \oplus\ZZ/2\{w\eta_8,j\eta_8\}.
 \end{split}
 \tag{D.2.22t7}
\]
The class $\tau$ is the Moore $\eta$ coextension in
$M(\ZZ/2,7)_R$, with $2\tau=R\eta_7^2$;
$w,j$ are the two free eight-sphere inclusions.
The stable stem and Moore identities used in these exact
sequences are the opened formulas
$\pi_1^S=\ZZ/2\{\eta\}$,
$\pi_2^S=\ZZ/2\{\eta^2\}$,
$\pi_3^S{}_{(2)}=\ZZ/8\{\nu\}$,
$\eta^3=4\nu$ and $2\widetilde\eta=i\eta^2$ in
\cite[\S2, (2.1)--(2.2), p.~149]{I_Oka1986}.

\medskip\noindent
\textbf{The exact $\eta$ composites of the free lifts.}
Set $J=M\cup_b e^8_{B_8}$.
Its cofiber sequence gives $\pi_9J=\ZZ/2\{\overline m\}$,
and $J/S^5$ is $C\eta$ with bottom cell in dimension six.
Let $e\in\pi_8J$ have top-cell degree two on $B_8$.
Its quotient is the top-degree-two class of $C\eta$.
The actual identity
$\eta_{C\eta}=i\nu q$, opened in
\cite[(8.4), p.~1279]{EH5AngeltveitRognes2005}, gives
\[
 e\eta_8=\overline m\quad\hbox{in }\pi_9J,
       \qquad e\eta_8=2s\quad\hbox{in }\pi_9L_8.
 \tag{D.2.22t8}
\]
Indeed its quotient is $2i\nu$, the nonzero image of
$\overline m$.  The free lift of $2C_8$ lies in the actual
$C\eta$ subcomplex on $r,C_8$ and has $\eta$ composite
$2r\nu=2z=0$.  The other two free composites are
$w\eta_8,j\eta_8$.

The degree-eight torsion of $L_8$ is
$\ZZ/4\{n=i_5\nu_5\}\oplus\ZZ/2\{\ell=R\eta_7\}$;
its free Hurewicz lattice has basis
$W,J_2,2J_1,2(D_8-J_1)$.
The nine-cells supply all four integral boundary directions.
The zero-boundary combinations $N$ and
$K_\Sigma=K_2+K_3+K_4$ kill $n,\ell$.
On taking $\eta$ composites, $n\eta=0$,
$\ell\eta=2\tau$, and \textup{(D.2.22t8)} gives $2s$.
Torsion additions to the integral attachments contribute only
$2\tau$, already killed.  Thus the image of
\textup{(D.2.22t7)} in $\pi_9L_9$ is exactly
\[
 V=(\ZZ/2)^5\{s,z,\tau,w\eta_8,j\eta_8\}.
 \tag{D.2.22t9}
\]
The quotient $\pi_9L_9/V$ is free of rank four, the kernel
of the full integral nine-cell attaching matrix.  In particular
$V$ is the entire torsion subgroup.  Its complete detectors are
\[
\begin{array}{c|c|c}
 &\operatorname{Sq}_*^2&\operatorname{Sq}_*^4\\ \hline
 s&0&D\\
 z&0&r\\
 \tau&J_1+J_3&0\\
 w\eta_8&W&0\\
 j\eta_8&J_2&0.
\end{array}
 \tag{D.2.22t10}
\]
The five vectors are independent.  Thus specifying these
coordinates leaves no unmeasured torsion component.

\medskip\noindent
\textbf{The free lattice, retaining the even-$\nu$ coefficients.}
Replace $N$ by $N'=N-K_\Sigma$, rescale $n$ by an odd
unit, and change $\ell$ by an even multiple of $n$.
We may then take the attachments of $N'$ and
$Q=K_\Sigma$ to be $n$ and $\ell$, respectively.
The order-eight homology generator changes only by an
order-two class.  The other two zero-boundary nine-cycles
have attachments
\[
 K_1:\ 2a n,
       \qquad K_5:\ 2b n,
       \qquad a,b\in\{0,1\}.
 \tag{D.2.22t11}
\]
These two bits are retained.  A mod-two spherical representative
alone would not set either bit to zero: its integral Hurewicz
class could differ from $K_i$ by $2N'$.

Choose integral order-two representatives $K_1,K_5,Q$ by
the half-boundaries of $B_3,B_7,B_4$, respectively, and use
the order-eight generator $N'$ supplied by the integral
configuration-chain calculation above.
The free Hurewicz lattice of $\pi_9L_9$ then has a basis
$p,q,x,y$ satisfying
\[
 \begin{aligned}
 h(p)&=4N',& h(q)&=2Q,\\
 h(x)&=K_1-2aN',&h(y)&=K_5-2bN'.
 \end{aligned}
 \tag{D.2.22t12}
\]
Normalize $p,q$ so that their second- and fourth-square
detectors are zero.  To justify this normalization, the Adem
relation $\operatorname{Sq}^2\operatorname{Sq}^2
=\operatorname{Sq}^3\operatorname{Sq}^1$ forces either
second-square vector into
$\langle W,J_2,J_1+J_3\rangle$.
For $p$ the top third square is zero; for $q$ it is $R$,
whose Bockstein is zero.  Either fourth-square vector lies
in $\langle D,r\rangle$.
The five entries of \textup{(D.2.22t10)} remove precisely
these five coordinates.  Even multiples of all attaching maps,
including $2x,2y$, have zero primary detectors by the
degree-two mapping-cone argument above.

\medskip\noindent
\textbf{All twelve ten-cell relations.}
The exact integral normalization is controlled by
$H_9E=\ZZ/8\{N'\}\oplus\ZZ/2\{K_1,K_5,Q\}$.
A first-Bockstein cycle with zero third Bockstein can be
made an exact zero-boundary ten-cell: its remaining $N'$
coefficient is divisible by sixteen and is removed by twice
an eight-$N'$ cell, while its order-two coefficients are
removed by even multiples of the corresponding degree-two
cells.  These changes preserve all primary detectors.
Since \emph{both} $B_5$ and $B_6$ have third Bockstein $N$,
they may each be normalized to boundary $8N'$.
Odd units do not change the ensuing two-primary relations.

The zero-boundary cells $Z_2,B_2,T_{10}$ give
\[
 w\eta_8=0,
       \qquad j\eta_8=0,
       \qquad w\eta_8+s+z=0.
 \tag{D.2.22t13}
\]
The two eight-$N'$ cells, using the normalized $p$, give
\[
 B_5:\ 2p+\tau+s=0,
       \qquad B_6:\ 2p+j\eta_8=0.
 \tag{D.2.22t14}
\]
Consequently $2p=0$ and $\tau=z=s$.
The zero-boundary cycle $B_1+B_3$ repeats $j\eta_8$.
The classes $U_2,P_r,B_8$ have zero detectors and impose
no new relation, by completeness of \textup{(D.2.22t10)}.
The third-Bockstein-zero combination $B_5+B_6$ repeats
$\tau+s+j\eta_8$.  Exactly one nonzero element of $V$
survives, represented by $s$.

The remaining integral cells have boundaries
$2Q,2K_1,2K_5$, which in \textup{(D.2.22t12)} are
$q,2x+ap,2y+bp$.  Their complete relations are
\[
 \begin{array}{ll}
 B_4:&q+\tau+s+j\eta_8=0,\\
 B_3:&2x+ap=0,\\
 B_7:&2y+bp=0.
 \end{array}
 \tag{D.2.22t15}
\]
Here $q$ was normalized above, and $2x,2y$ have zero
primary detectors; $B_1$ differs from $B_3$ by the already
killed $j\eta_8$.  All twelve ten-cells and all independent
integral boundary directions have now been used.
The resulting full presentation is
\[
 \pi_9E=
 \langle p,x,y,s\mid
       2p=2s=0,\quad2x=ap,\quad2y=bp\rangle.
 \tag{D.2.22t16}
\]
In particular, independently of the two undetermined bits,
\[
 4\pi_9T=0,
 \qquad
 \ker\bigl(h_{\ZZ}:\pi_9E\to H_9E\bigr)
       =\ZZ/2\{s\}\ne0.
 \tag{D.2.22t17}
\]
Its Hurewicz image is given exactly by
\textup{(D.2.22t12)} and has order eight.
The group is $(\ZZ/2)^4$ if $a=b=0$, and
$\ZZ/4\oplus(\ZZ/2)^2$ otherwise.
These alternatives are retained rather than selected by
an abstract group-order argument.
The previously proved injectivity of $H_9E\to H_9F$
and the splitting $F\to T$ identify the same kernel for
the integral Hurewicz map on $\pi_9T$.

Finally the $r$ sphere maps to an odd multiple of the
class $\rho\in\pi_6T=\ZZ/4$ of D.2.18.
Since $z=r\nu_6$ and $z=s$ in the final group, the
surviving element has the specified actual representative
\[
 s=\rho\nu_6\ne0.
 \tag{D.2.22t18}
\]
The same eight- and nine-cell exact sequences kill $n,\ell$
and the free directions $2J_1,2(D_8-J_1)$, while leaving
the $W,J_2$ directions modulo two.
Thus the actual group $\pi_8E=\pi_8T\cong(\ZZ/2)^2$ is
represented by those two eight-spheres.  Their $\eta$
composites are precisely the two classes killed in
\textup{(D.2.22t13)}.  Hence
\[
 \pi_8T\circ\eta_8=0.
 \tag{D.2.22t19}
\]
This excludes that particular proposed change of a relative
eight-cell filling.  It does not exhaust every possible change
of the preceding map and does not evaluate
$\lambda(\chi,\chi)$.
